\documentclass[11pt]{article}
\usepackage[margin=1in]{geometry}
\usepackage{amsmath}
\title{Cryogenic Kinetics Models}
\begin{document}
\maketitle

This note collects the simple temperature-dependent kinetics models used in the cryogenic extension of the project.

\section*{1. Defect migration}
A first-order Arrhenius model for defect diffusivity is
\[
D_{\mathrm{def}}(T)=D_0 \exp\!\left(-\frac{E_m}{k_B T}\right).
\]
Taking the logarithm gives
\[
\ln D_{\mathrm{def}} = \ln D_0 - \frac{E_m}{k_B T}.
\]
This predicts an approximately straight line on an Arrhenius plot of $\ln D_{\mathrm{def}}$ versus $1/T$.

\section*{2. Trap emission}
For thermally activated emission from a defect level, use
\[
e(T)=e_0 \exp\!\left(-\frac{E_a}{k_B T}\right).
\]
The corresponding emission time constant is
\[
\tau_{\mathrm{emit}}(T)=\frac{1}{e(T)}
\]
and, if $e_0 = 1/\tau_0$,
\[
\tau_{\mathrm{emit}}(T)=\tau_0 \exp\!\left(\frac{E_a}{k_B T}\right).
\]
Thus the trap-emission timescale grows rapidly as temperature decreases.

\section*{3. Annealing}
A first-order annealing-rate model is
\[
R_{\mathrm{ann}}(T)=R_0 \exp\!\left(-\frac{E_{\mathrm{ann}}}{k_B T}\right).
\]
A simple annealing timescale can be defined as
\[
\tau_{\mathrm{ann}}(T)=\frac{1}{R_{\mathrm{ann}}(T)}.
\]
This makes clear that defect recovery can become extremely slow in the cryogenic regime.

\section*{4. Interpretation}
These models are not meant to represent one unique microscopic defect species. Instead, they are first-order phenomenological models used to illustrate how thermally activated kinetics collapse as the system moves toward low temperature.

\end{document}
